%
%
%
%
%
\definecolor{skI}{HTML}{0072B2}\definecolor{skII}{HTML}{D55E00}\definecolor{skIII}{HTML}{009E73}\definecolor{skgrau}{HTML}{5F6368}
\tikzset{sk/duenn/.style={line width=0.5pt}, sk/dick/.style={line width=1.8pt}}
\ifnum\pdfstrcmp{\teilnr}{I}=0 \tikzset{sk/I/.style={sk/dick}}\else\tikzset{sk/I/.style={sk/duenn}}\fi
\ifnum\pdfstrcmp{\teilnr}{II}=0 \tikzset{sk/II/.style={sk/dick}}\else\tikzset{sk/II/.style={sk/duenn}}\fi
\ifnum\pdfstrcmp{\teilnr}{III}=0 \tikzset{sk/III/.style={sk/dick}}\else\tikzset{sk/III/.style={sk/duenn}}\fi
\ifnum\pdfstrcmp{\teilnr}{0}=0 \tikzset{sk/band/.style={sk/dick}}\else\tikzset{sk/band/.style={sk/duenn}}\fi
\providecommand{\skp}[1]{\par\hangindent=0.75em\hangafter=1\noindent\makebox[0.75em][l]{\textbullet}#1}
%
\providecommand{\skspalte}[1]{\parbox[t][4.8cm][t]{4.5cm}{\raggedright\hyphenpenalty=10000\exhyphenpenalty=10000 #1}}
\providecommand{\skdies}[1]{\ifnum\pdfstrcmp{\teilnr}{#1}=0 \\[1pt]\textbf{this document}\fi}
%
%
\par\medskip\noindent\begin{minipage}{\textwidth}
\centering
\begin{tikzpicture}[font=\footnotesize,
  spalte/.style={rectangle, rounded corners=2pt, inner sep=5pt, anchor=north west},
  einzel/.style={rectangle, rounded corners=2pt, align=center, text width=4.0cm, inner sep=4pt, minimum height=0.95cm, anchor=north,
    execute at begin node={\hyphenpenalty=10000\exhyphenpenalty=10000}},
  pfeil/.style={-{Stealth[length=2mm]}, draw=skgrau, line width=0.5pt},
  quer/.style={{Stealth[length=1.6mm]}-{Stealth[length=1.6mm]}, draw=skgrau, line width=0.5pt, dashed}]
\node[spalte, draw=skI, fill=skI!5] (sI) at (0, 0) {\skspalte{\textbf{\color{skI}Part I: triples}\par\smallskip
  {\color{skgrau}Can $n - 1$, $n$, $n + 1$ all be powerful? (Erd\H{o}s)}\par\smallskip
  \skp{the middle is a Pell number $T_k(m)$ at a point $(m, k)$}
  \skp{a height inequality, two wings of kernels}
  \skp{witnesses: primes dividing $T_k$ exactly once}
  \skp{no triple with a middle below $\Z{ohnetripelsicher}$}
  \skp{\Z{rgpaare} building blocks, \Z{rgoffen} open}\par}};
\node[spalte, draw=skII, fill=skII!5] (sII) at (5.05, 0) {\skspalte{\textbf{\color{skII}Part II: pairs}\par\smallskip
  {\color{skgrau}How often are $n$ and $n + 1$ both powerful? (Erd\H{o}s)}\par\smallskip
  \skp{Walker's families; all pairs up to $2^{\Z{paare52exp}}$}
  \skp{$\liminf N_1(x)/\log x \ge \Z{csechs}$}
  \skp{distance two: a sum over towers, a lower bound at every height}
  \skp{upper bounds reduce to the kernels with $m \mid U_1(m)$}\par}};
\node[spalte, draw=skIII, fill=skIII!5] (sIII) at (10.1, 0) {\skspalte{\textbf{\color{skIII}Part III: observations}\par\smallskip
  {\color{skgrau}What else did the study of triples and pairs show?}\par\smallskip
  \skp{the gaps between consecutive powerful numbers}
  \skp{powerful values of cyclotomic polynomials}
  \skp{tests of the independence assumptions, with their power}\par}};
\path let \p1 = (sI.south), \p2 = (sII.south), \p3 = (sIII.south) in
  node[rectangle, rounded corners=2pt, draw=skgrau, fill=skgrau!6, align=center, text width=14.67cm, inner sep=4pt, anchor=north west]
  (rahmen) at (0pt, {min(\y1, \y2, \y3) - 6pt}) {\mbox{framework introduction} \enspace$\cdot$\enspace \mbox{results at a glance} \enspace$\cdot$\enspace
  \mbox{known open problems} \enspace$\cdot$\enspace \mbox{Problems A to \problemeletzter} \enspace$\cdot$\enspace \mbox{data and code availability}};
\node[anchor=south west, font=\footnotesize\bfseries, inner sep=1pt] (bandkopf) at (0, 0.12) {The collected volume};
\begin{scope}[on background layer]
\node[rectangle, rounded corners=4pt, draw=black, sk/band, fit=(sI)(sII)(sIII)(rahmen)(bandkopf), inner sep=6pt] (band) {};
\end{scope}
\node[einzel, draw=skI, fill=skI!5, sk/I] (eI) at ($(sI.north |- band.south) + (0, -0.75)$) {\textbf{\color{skI}Part I} as a paper of its own\skdies{I}};
\node[einzel, draw=skII, fill=skII!5, sk/II] (eII) at ($(sII.north |- band.south) + (0, -0.75)$) {\textbf{\color{skII}Part II} as a paper of its own\skdies{II}};
\node[einzel, draw=skIII, fill=skIII!5, sk/III] (eIII) at ($(sIII.north |- band.south) + (0, -0.75)$) {\textbf{\color{skIII}Part III} as a paper of its own\skdies{III}};
\foreach \x in {I, II, III} \draw[pfeil] (s\x.south |- band.south) -- (e\x.north);
\draw[quer] (eI.east) -- (eII.west);
\draw[quer] (eII.east) -- (eIII.west);
\path let \p1 = (eI.south), \p2 = (eII.south), \p3 = (eIII.south) in
  node[font=\scriptsize, text=skgrau, anchor=north, text width=13.5cm, align=center] at ($(eII.south |- 0pt, {min(\y1, \y2, \y3)}) + (0, -0.08)$)
  {each part also stands as a paper of its own, with the same text and numbering; dashed: the parts refer to each other};
\ifnum\pdfstrcmp{\teilnr}{0}=0 \node[font=\scriptsize\bfseries, anchor=south east, inner sep=1pt] at (band.north east) {this document};\fi
\end{tikzpicture}
\captionfigur{\textbf{What this research did.} It followed two questions of Erd\H{o}s on consecutive powerful numbers, on triples and on pairs,
through Pell equations, proved or computed what it could reach, and recorded the rest as observations and open problems. The collected volume
contains the three parts and the framework that ties them together; each part also stands as a paper of its own, with the same text and numbering,
and the parts refer to each other. The heavier frame marks the document in hand.}
\end{minipage}\par\medskip
